% Discussion and Conclusions. Result statements map to claims.tex; limitations follow the plan (Section 1) and the
% related-work passes.
\section{Discussion}\label{sec:discussion}

\paragraph{Efficiency does not show that a prior-weighted search keeps what its prior disfavours.} In every primary setting
some weight of the prior gave a demonstrable gain; at $\beta=0.5$, all \R{benefcount/primary/0.5} settings did, and none was
safe under the registered criterion. A practitioner who chose the weight by speed alone would therefore have chosen weights
for which a loss of a fifth or more of the baseline's recall of the good materials the prior disfavours could not be
excluded. The settings differed in how soon this happened. For the elemental-modulus heuristic on the shear modulus, the gain
was large, the loss of $\Sdis$ was shown to be below 20\% of the baseline at every weight up to $\beta=0.25$, and no loss was
detected even at $\beta=0.5$; for the shear-modulus form of the same heuristic applied to the
specific modulus, a large share of $\Sdis$ was lost at the smallest weight, for a gain that was not demonstrable. Matching the
heuristic to the objective raised its gain on Y2, to \R{frontier/primary/Y2:P_phys_matched_Y2/0.25/gain}
\ci{frontier/primary/Y2:P_phys_matched_Y2/0.25/gain} acquisitions at $\beta=0.25$, and H1 detected no loss for it after Holm
correction (although its unadjusted intervals exclude zero), but no weight was safe: the bounds on its loss could not exclude 20\% of the baseline recall, and at $\beta=0.05$, $0.1$ and $0.75$ even the
unadjusted one-sided 95\% bound showed a loss of more than 20\%. Absence of a detected loss and a demonstrably small loss are different
statements, and only the second supports using a weight.

\paragraph{Redirection, not lower yield.} In \R{n5up/primary/0.5} of the 6 settings the weighted search acquired more top-5\%
materials in total than the no-prior search at $\beta=0.5$, while its estimated recall of $\Sdis$ was lower, a loss detected by H1 in
\R{n5up_h1/primary/0.5} of them (Section~\ref{sec:where}). For most priors here,
then, the cost of the weight is not fewer good materials but different ones: those the prior already favours. Whether that cost
matters depends on the purpose of the search. It does not matter for a search that only needs some good materials; it does for one
meant to map which good materials exist, or to find those that an incomplete heuristic or a biased history undervalues. The
exception, the shear-modulus form of the elemental heuristic applied to the specific modulus, lost coverage without finding more
good materials.

\paragraph{What is built into the coverage measure, and what is informative.} Because $\Sdis$ is the half of the good materials
that the prior ranks lowest, a search that follows the prior closely must lose it: at $\beta=1$ its recall was zero in every cell, as the
construction leads one to expect, and the loss there carries little information. What the measure adds is how the loss develops at the
weights a practitioner would use: the weight at which it becomes detectable or exceeds the safety margin, how large it is when the
gain first becomes demonstrable, and how these differ between settings. On those questions the six settings differed widely,
from an estimated loss of \R{h2detail/primary/Y2:P_phys_fixed/0.02/relative_loss_pct}\% of the baseline at $\beta=0.02$ (the
Y2 elemental-modulus setting) to no loss detected by H1 even at $\beta=0.5$ (the Y1 elemental-modulus and Y2
elemental-modulus-matched settings). $N_5$, whose definition does not involve the prior, gives the total yield beside it.

\paragraph{What coverage measures.} $\Sdis$ asks whether weighting a prior costs the search the good materials that this prior
ranks lowest. Because it is built from the prior under test, it reflects how the prior ranks the best materials as much as
how the search behaves: a prior that ranks almost all good materials highly has an $\Sdis$ that lies mostly above its 80th
percentile, where a prior-weighted search still looks, so $\Sdis$ can be kept while the few good materials below that
percentile are lost. The absolute blind spot $\Sblind$ separates these effects. For the
elemental-modulus heuristic on Y1, no loss of $\Sdis$ was detected at $\beta=0.5$, but almost all of its small blind spot was
lost (post hoc). In this benchmark, that is the clearest sign of where a strong fixed weight takes the search: away from the good
materials the prior ranks low, even when the prior is good enough that few such materials exist. Neither set
measures diversity or global exploration, and other definitions of the materials a search should not miss would give other
numbers.

\paragraph{The two priors differ in information as well as in form.} The historical prior scores a candidate by its distance
to the best materials of a labelled history of about a thousand structures. It is a nearest-neighbour heuristic by construction,
it inherits whatever the history over-represents\cite{jia2019anthropogenic}, and such similarity-based scores are known to reach
poorly beyond the chemistries they were built on\cite{meredig2018loco}. A coverage loss is what such a construction is likely to
produce; what the benchmark adds is its size and how it grows with $\beta$. The
elemental-modulus heuristic uses no labels at all. Differences between the two priors are differences between these complete
strategies, not between prior forms at equal information, and they should not be read as showing that heuristics beat data or
the reverse.

\paragraph{Fixed and decaying weights.} The weight studied here does not decay. Prior-guided methods such as $\pi$BO let the
prior's influence fade as observations accumulate, so that the search recovers from a misleading
prior\cite{hvarfner2022pibo,souza2021bopro}; recent variants decay a prior placed in the Gaussian-process mean\cite{qu2026dynmeanbo}
or prior-based preferences in the acquisition function\cite{fehring2025dynabo}. This study cannot say whether a decaying weight avoids the coverage loss: a fixed and a
decaying weight are not on a common scale, and losses appeared within the first 10 acquisitions, when a decaying prior's influence
is strongest. Measuring the coverage of decaying schemes on the same benchmark is a direct next step.

\paragraph{Alignment.} Across the six settings, the history-side correlation between prior and objective went together with the
gain at $\beta=0.5$ (Kendall $\tau=\R{summary/tau}$, exploratory, pooling both objectives; with six settings this is compatible with no
association). Six settings that share objectives,
prior constructions and a pool cannot establish that alignment explains the frontier. The observation suggests a hypothesis that
can be tested before any search is run, that history-side alignment predicts a prior's gain; it says nothing about coverage, and
the most aligned prior on Y2 had no safe weight. Testing it as a prediction on other pools and properties would turn it into
evidence.

\paragraph{Limitations.} Several follow from the design and are stated in the registration.
\begin{itemize}\setlength{\itemsep}{1pt}
\item \emph{Scope.} One database snapshot, metallic multi-element structures, DFT elastic properties, a noiseless lookup oracle
over already-relaxed structures, one surrogate, one acquisition function, batch size 1 and a budget of 10\,+\,40 evaluations. The
intervals describe variation over partitions of one fixed pool, whose discovery sides share about 70\% of their materials, not
over pools or databases. The chemical-system splits hold out whole chemical systems, but 24--31\% of their discovery materials
have a proper sub-system in history (for example, a binary of a held-out ternary), so their discovery materials are not far from
what the history covers.
\item \emph{Curation.} The primary universe keeps only structures whose elements pass the screen of the elemental prior's own
reference data, which removed 24.6\% of the structures; this may favour that prior. \R{bench/snap/primary_theoretical_pct}\% of the primary structures are hypothetical
(MP's \emph{theoretical} flag), and no thermodynamic-stability filter is applied.
\item \emph{Descriptors.} The descriptors include each candidate's DFT-relaxed density and volume, and Y2 divides by that density,
so the Y2 surrogate sees part of its objective.
\item \emph{Surrogate.} The Gaussian process uses one isotropic length-scale, bounded below at 3.0; the fitted length-scale was
at that bound in \R{healthrange/primary/lb_min}--\R{healthrange/primary/lb_max}\% of steps per arm, where shorter-range structure
could not be represented. In pool-based materials benchmarks, anisotropic
Gaussian processes and random forests have outperformed isotropic ones\cite{liang2021benchmarking}, so the no-prior baseline here
may be weaker than a tuned one. Because the prior arms mix their ranks with the same surrogate, the direction in which a
stronger surrogate would change the gains, and the baseline recall of $\Sdis$, is not known.
\item \emph{Outcomes.} The efficiency outcome is a restricted mean; where the shares of runs without a hit differ greatly, as for
the prior alone on Y2, it must be read with them. $\Sdis$ has some loss at high weight built in, and its recall levels are not
comparable across priors.
\item \emph{Decision rule.} The 20\% margin is a judgement. The registered rule controls false declarations of safety within a
cell, not the joint beneficial-and-safe call, and nothing is adjusted across cells; with both bounds adjusted (post hoc), the
Y2/historical-matched setting no longer qualified, since its weight rested on the unadjusted beneficial bound. The weight $\beta$
is not comparable to the decaying weights of $\pi$BO-type methods, and no decaying arm was run. The qualifying weights
also changed when Pm and Tc were removed or the reference screen was dropped.
\item \emph{Exploratory history.} The benchmark, outcomes and hypotheses were developed with access to outcome labels, including
the discovery sides of three splits of each type, and engine choices were informed by searches on history sides that cover most
of every discovery pool. The analysis plan was registered, and amended once, before any model-driven search was run for
results. The registered analysis guards the tests, not the choice of what to test; the previously unanalysed
splits r3--r9 are a robustness check within the same benchmark, not a replication.
\end{itemize}

\paragraph{Implications.} Prior-guided search is spreading to priors whose reliability is hard to judge in advance, including
those supplied by language models\cite{chen2026evidencegated} or language-model-guided search\cite{liu2024llambo}. Our results argue for a concrete reporting practice: beside any claim of faster discovery, report, at the weight actually
used, the total number of good materials found ($N_5$ here) and the recall of those that the prior disfavours ($\Sdis$ here), each
against the no-prior search on the same initial designs. Neither number alone shows whether a prior redirects the search or
narrows it.
The benchmark, its run records and its analysis code are available for testing other priors, weights and surrogates against the
same coverage endpoints.

\section{Conclusions}\label{sec:conclusions}
In a benchmark with a registered analysis on Materials Project elasticity data, weighting a scientific prior in Bayesian search
traded efficiency, at sufficient weight, against the recall of good materials that the prior disfavours, and the terms of that trade differed between
prior--objective settings. A weight that was both beneficial and safe under the registered criterion existed in three of six
settings (two under a post hoc joint error control), was always small, and depended on the safety margin and on the curation of
the benchmark. No weight qualified in every setting, and strong fixed weights lost the prior's blind spot even where the prior
was well aligned with the objective (post hoc). In most settings the weighted search found more good materials in total but not the same ones. Faster discovery alone is
therefore not evidence that a prior-weighted search keeps the materials its prior disfavours; that has to be measured. These
results are relative to one simple surrogate, a Gaussian process with a single isotropic Mat\'ern-5/2 length-scale, and expected
improvement; a richer surrogate, for example with a length-scale per descriptor, may change them.
